% physics_wall_crossing_mc.tex
% Physics note: Wall-crossing as gauge equivalence of Maurer-Cartan elements
% Raeez Lorgat, 2026

\documentclass[11pt]{article}
\usepackage[T1]{fontenc}
\usepackage[utf8]{inputenc}
\usepackage{lmodern}
\frenchspacing

\usepackage{amsmath,amssymb,amsthm}
\usepackage{mathtools}
\usepackage{enumitem}
\usepackage[margin=1.2in]{geometry}
\usepackage{hyperref}

% Theorem environments
\newtheorem{theorem}{Theorem}[section]
\newtheorem{proposition}[theorem]{Proposition}
\newtheorem{lemma}[theorem]{Lemma}
\newtheorem{corollary}[theorem]{Corollary}
\newtheorem{conjecture}[theorem]{Conjecture}
\theoremstyle{definition}
\newtheorem{definition}[theorem]{Definition}
\newtheorem{construction}[theorem]{Construction}
\newtheorem{example}[theorem]{Example}
\theoremstyle{remark}
\newtheorem{remark}[theorem]{Remark}
\newtheorem{warning}[theorem]{Warning}

% Macros (subset from main.tex)
\newcommand{\frakg}{\mathfrak{g}}
\newcommand{\frakh}{\mathfrak{h}}
\newcommand{\frakn}{\mathfrak{n}}
\newcommand{\Ainf}{A_\infty}
\newcommand{\Linf}{L_\infty}
\newcommand{\Mod}{\mathrm{Mod}}
\newcommand{\Rep}{\mathrm{Rep}}
\newcommand{\Hom}{\mathrm{Hom}}
\newcommand{\End}{\mathrm{End}}
\newcommand{\id}{\mathrm{id}}
\newcommand{\ad}{\mathrm{ad}}
\newcommand{\MC}{\mathrm{MC}}
\newcommand{\CY}{\mathrm{CY}}
\newcommand{\HH}{\mathrm{HH}}
\newcommand{\Tr}{\mathrm{Tr}}
\newcommand{\DT}{\mathrm{DT}}
\newcommand{\Coh}{\mathrm{Coh}}
\newcommand{\Fuk}{\mathrm{Fuk}}
\newcommand{\Ran}{\mathrm{Ran}}
\newcommand{\Sym}{\mathrm{Sym}}
\newcommand{\cA}{\mathcal{A}}
\newcommand{\cB}{\mathcal{B}}
\newcommand{\cC}{\mathcal{C}}
\newcommand{\cM}{\mathcal{M}}
\newcommand{\cO}{\mathcal{O}}
\newcommand{\cR}{\mathcal{R}}
\newcommand{\cE}{\mathcal{E}}
\newcommand{\cN}{\mathcal{N}}
\newcommand{\cW}{\mathcal{W}}
\newcommand{\cZ}{\mathcal{Z}}
\DeclareMathOperator{\mult}{mult}
\DeclareMathOperator{\Aut}{Aut}
\DeclareMathOperator{\Stab}{Stab}
\DeclareMathOperator{\Arg}{Arg}

\title{Wall-Crossing as Gauge Equivalence of Maurer--Cartan Elements\\[6pt]
\large Physics Note for \emph{Calabi--Yau Quantum Groups}}
\author{Raeez Lorgat}
\date{April 2026}

\begin{document}
\maketitle

\begin{abstract}
We explain how the Kontsevich--Soibelman wall-crossing formula for BPS
degeneracies of a Calabi--Yau threefold $X$ admits a natural interpretation
in the quantum vertex chiral group framework: wall-crossing is gauge
equivalence in the Maurer--Cartan moduli space of the modular convolution
algebra $\frakg^{\mathrm{mod}}_A$ associated to the chiral algebra
$A = \Phi(\cC)$. The KS product formula is identified with the gauge
invariance of the bar-complex Euler product (the denominator identity of
the BKM superalgebra $\frakg_X$), the attractor mechanism corresponds to a
canonical gauge for the MC element determined by the split BKM root datum,
and the discontinuity of the shadow obstruction tower truncations $\Theta^{\leq r}_A$
at walls of marginal stability is the obstruction to gauge invariance at
finite order.
\end{abstract}

\tableofcontents

% ======================================================================
\section{Introduction and motivation}
\label{sec:wc-intro}
% ======================================================================

Let $X$ be a Calabi--Yau threefold and $\cC = D^b(\Coh(X))$ (or $\Fuk(X)$)
the associated CY$_3$ category. The BPS spectrum of $X$ is the function
\[
 \Omega \colon \Gamma \longrightarrow \mathbb{Z}, \qquad
 \gamma \longmapsto \Omega(\gamma; t),
\]
where $\Gamma = H_{\mathrm{even}}(X, \mathbb{Z})$ (or $K_0(\cC)$) is the
charge lattice and $t$ denotes the complexified K\"ahler moduli. The
celebrated result of Kontsevich--Soibelman \cite{KS08} is that the
$\Omega(\gamma; t)$ are \emph{not} invariant under variation of $t$:
they jump at \emph{walls of marginal stability}, codimension-one loci in
moduli space where BPS bound states form or decay. Nevertheless, the KS
wall-crossing formula organizes these jumps into a single invariant datum
: a product of symplectomorphisms of a formal torus algebra.

The purpose of this note is to explain, in the language of the quantum
vertex chiral group programme (Volumes~I--III), \emph{what} this invariant
datum is: it is a gauge equivalence class in the Maurer--Cartan moduli space
of the modular convolution algebra $\frakg^{\mathrm{mod}}_A$. The principal
claims are:

\begin{enumerate}[label=(\arabic*)]
\item The MC element $\Theta_A \in \MC(\frakg^{\mathrm{mod}}_A)$ depends on
 the CY moduli $t$. As $t$ varies, $\Theta_A(t)$ traces a path in the MC
 moduli space.

\item Crossing a wall of marginal stability is a gauge transformation
 $\Theta_A \mapsto e^{\alpha} \cdot \Theta_A$ where $\alpha$ encodes the
 BPS states that bind or decay.

\item The KS wall-crossing formula
 $\prod_{\gamma \text{ on wall}} (1 - x^\gamma)^{\Omega(\gamma)} = \text{const}$
 is precisely the statement that the bar-complex Euler product
 (= denominator identity of the BKM superalgebra $\frakg_X$) is
 gauge-invariant.

\item The attractor mechanism places $\Theta_A$ in a canonical
 (``split'') gauge at the attractor point; away from the attractor, gauge
 transformations reshuffle root multiplicities.

\item The shadow obstruction tower truncation $\Theta^{\leq r}_A$ may jump
 discontinuously at walls (finite-order truncations are \emph{not}
 gauge-invariant), but the full $\Theta_A$ is.
\end{enumerate}

This note is expository. Precise statements at the level of the main text
will appear in Part~\ref{part:frontier} (Connections) once the full apparatus of
Chapters~\ref*{ch:modular-trace}--\ref*{ch:bar-cobar-bridge} is in place.

% ======================================================================
\section{Recollections: MC elements and gauge equivalence}
\label{sec:wc-mc-recollections}
% ======================================================================

\subsection{The dg Lie algebra and its MC equation}
\label{ssec:dgla-mc}

Let $(\frakg, d, [\cdot, \cdot])$ be a differential graded Lie algebra
(dgla) over a field $k$ of characteristic zero. The \emph{Maurer--Cartan
equation} in $\frakg^1$ (degree-1 elements, cohomological convention) is
\begin{equation}\label{eq:mc-equation}
 d\Theta + \tfrac{1}{2}[\Theta, \Theta] = 0.
\end{equation}
We write $\MC(\frakg) = \{ \Theta \in \frakg^1 \mid
d\Theta + \frac{1}{2}[\Theta, \Theta] = 0 \}$ for the set of MC elements.

More generally, if $\frakg$ is a curved $\Linf$-algebra with operations
$\ell_0, \ell_1, \ell_2, \ldots$, the MC equation becomes
\begin{equation}\label{eq:mc-linf}
 \sum_{n=0}^{\infty} \frac{1}{n!} \ell_n(\Theta, \ldots, \Theta) = 0.
\end{equation}
The modular convolution algebra $\frakg^{\mathrm{mod}}_A$ of Volume~I is
a curved $\Linf$-algebra of precisely this type, where the curvature
$\ell_0$ encodes the genus-zero data (the Weyl vector contribution) and
the higher operations $\ell_n$ encode $n$-point interactions on moduli
spaces of curves.

\subsection{Gauge equivalence}
\label{ssec:gauge-equiv}

Two MC elements $\Theta, \Theta'$ are \emph{gauge equivalent} if they are
connected by the action of the gauge group $\exp(\frakg^0)$. Concretely,
for $\alpha \in \frakg^0$:
\begin{equation}\label{eq:gauge-action}
 e^{\alpha} \cdot \Theta
 = \Theta + \sum_{n=0}^{\infty} \frac{(\ad \alpha)^n}{(n+1)!}
 \bigl( d\alpha + [\alpha, \Theta] \bigr).
\end{equation}
This is the Baker--Campbell--Hausdorff gauge action. The MC
\emph{moduli space} is the quotient
\[
 \overline{\MC}(\frakg) = \MC(\frakg) / \exp(\frakg^0),
\]
and it is the deformation-theoretic invariant: gauge-equivalent MC elements
give equivalent deformation problems.

In the $\Linf$ setting, the gauge action generalizes: for $\alpha \in
\frakg^0$, the homotopy transfer produces a path
$\Theta_s$, $s \in [0,1]$, in $\MC(\frakg)$ with
\begin{equation}\label{eq:linf-gauge}
 \frac{d\Theta_s}{ds}
 = \sum_{n=0}^{\infty} \frac{1}{n!}
 \ell_{n+1}(\alpha, \Theta_s, \ldots, \Theta_s).
\end{equation}
The endpoint $\Theta_1 = e^{\alpha} \cdot \Theta_0$ is the gauge-transformed
MC element.

\subsection{The modular convolution algebra}
\label{ssec:modular-conv}

For a chiral algebra $A$ (Vol.~I), the modular convolution algebra
$\frakg^{\mathrm{mod}}_A$ is the (curved) $\Linf$-algebra whose
\begin{itemize}
\item underlying graded vector space is
 $\prod_{g \geq 0} H^\bullet(\overline{\cM}_{g,n}, \cE^{(g)}_A)$
 (sections of the chiral insertion sheaf over moduli of curves),
\item $\Linf$-operations $\ell_n$ come from degeneration of curves
 (boundary strata of $\overline{\cM}_{g,n}$),
\item curvature $\ell_0 \in \frakg^{\mathrm{mod},2}_A$ is the genus-zero
 obstruction.
\end{itemize}
The universal MC element $\Theta_A \in \MC(\frakg^{\mathrm{mod}}_A)$ encodes
the full modular data of $A$: its degree-$r$ projections
$\Theta^{\leq r}_A$ form the shadow obstruction tower of Volume~I.

% ======================================================================
\section{The MC element of a CY3 category}
\label{sec:wc-mc-cy3}
% ======================================================================

\subsection{Setup}
\label{ssec:cy3-setup}

Let $X$ be a CY3, $\cC = D^b(\Coh(X))$, and $A = A_\cC = \Phi(\cC)$ the
quantum chiral algebra produced by the conjectural CY-to-chiral functor (CY-A;
proved for toric CY3, open for compact CY3 at $d = 3$).
The modular convolution algebra $\frakg^{\mathrm{mod}}_A$ is then a curved
$\Linf$-algebra whose
\begin{itemize}
\item charge lattice is $\Gamma = K_0(\cC) \cong H_{\mathrm{even}}(X,
 \mathbb{Z})$,
\item degree-1 elements are formal sums $\Theta = \sum_\gamma
 \Theta_\gamma \, x^\gamma$ over the charge lattice,
\item bracket is determined by the Euler form
 $\chi(\gamma, \gamma') = \sum_i (-1)^i \dim \mathrm{Ext}^i(E, F)$
 for $[E] = \gamma$, $[F] = \gamma'$.
\end{itemize}

The MC element $\Theta_A$ depends on the complexified K\"ahler moduli
$t \in \cM_K(X)$ through the stability condition:
\begin{equation}\label{eq:theta-t-dependence}
 \Theta_A(t) = \sum_{\gamma \in \Gamma_+}
 \Omega(\gamma; t) \cdot \Theta_\gamma^{\mathrm{prim}}(t) \, x^\gamma
\end{equation}
where $\Omega(\gamma; t)$ is the (rational) BPS index of charge $\gamma$
at moduli $t$, $\Theta_\gamma^{\mathrm{prim}}$ is the ``primitive''
MC contribution from a single BPS particle of charge $\gamma$, and
$\Gamma_+ \subset \Gamma$ is the positive cone determined by the stability
condition.

\begin{remark}[Dependence on stability condition]\label{rem:stability}
More precisely, $\Theta_A$ depends on a Bridgeland stability condition
$\sigma = (Z, \cA) \in \Stab(\cC)$, where $Z \colon \Gamma \to \mathbb{C}$
is the central charge and $\cA$ is the heart. The complexified K\"ahler
moduli $t$ determine $Z$ via the period map $t \mapsto Z_t(\gamma) =
-\int_\gamma e^{-(B + iJ)}$, where $B + iJ \in H^{1,1}(X, \mathbb{C})$.
The MC element thus traces a path in $\MC(\frakg^{\mathrm{mod}}_A)$
parameterized by $\Stab(\cC)$.
\end{remark}

\subsection{The bar-complex Euler product}
\label{ssec:bar-euler}

The bar complex $B(A)$ of the chiral algebra $A = A_\cC$ is a
factorization coalgebra on $\Ran(C)$. Its Euler characteristic (the
``denominator identity'' of the BKM superalgebra $\frakg_X$) is the
formal product
\begin{equation}\label{eq:euler-product}
 \Phi_X(x)
 = \prod_{\gamma \in \Gamma_+}
 (1 - x^\gamma)^{(-1)^{|\gamma|} \mult(\gamma)}
\end{equation}
where $\mult(\gamma)$ is the dimension of the root space
$(\frakg_X)_\gamma$ and $|\gamma|$ is the fermion parity. For
$X = K3 \times E$, this is the product formula for $\frac{1}{64}\Delta_5$
(Theorem~\ref*{thm:k3e-product} of Ch.~\ref*{ch:k3-times-e}).

The critical observation is:

\begin{proposition}[Euler product is gauge-invariant]
\label{prop:euler-gauge-inv}
The bar-complex Euler product $\Phi_X(x)$ depends only on the gauge
equivalence class $[\Theta_A] \in \overline{\MC}(\frakg^{\mathrm{mod}}_A)$,
not on the representative $\Theta_A$.
\end{proposition}

\begin{proof}[Proof sketch]
The Euler product is the graded character of the bar complex:
\[
 \Phi_X = \sum_\gamma (-1)^{|\gamma|}
 \dim H^\bullet(B(A), d_B)_\gamma \cdot x^\gamma.
\]
A gauge transformation $\Theta_A \mapsto e^\alpha \cdot \Theta_A$
induces a quasi-isomorphism of bar complexes
\[
 B(A; \Theta_A) \xrightarrow{\;\sim\;} B(A; e^\alpha \cdot \Theta_A)
\]
since the gauge transformation lifts to a chain map.
Quasi-isomorphic complexes have the same Euler characteristic.
\end{proof}

% ======================================================================
\section{Wall-crossing as gauge transformation}
\label{sec:wc-gauge}
% ======================================================================

\subsection{Walls of marginal stability}
\label{ssec:walls}

A \emph{wall of marginal stability} for the decomposition
$\gamma = \gamma_1 + \gamma_2$ is the real-codimension-one locus
\begin{equation}\label{eq:wall}
 \cW(\gamma_1, \gamma_2)
 = \bigl\{ t \in \cM_K(X)
 \,\big|\, \Arg Z_t(\gamma_1) = \Arg Z_t(\gamma_2) \bigr\}
 \subset \cM_K(X).
\end{equation}
At $t \in \cW(\gamma_1, \gamma_2)$, the BPS particles of charges
$\gamma_1$ and $\gamma_2$ are mutually BPS (their central charges are
aligned), and bound states of total charge $\gamma$ can form or decay.

As $t$ crosses $\cW(\gamma_1, \gamma_2)$, the BPS index $\Omega(\gamma; t)$
jumps. The Kontsevich--Soibelman wall-crossing formula constrains the jump.

\subsection{The KS formula via ordered products}
\label{ssec:ks-formula}

For each $\gamma \in \Gamma$ with $\Omega(\gamma) \neq 0$, define the
\emph{KS symplectomorphism} of the formal torus algebra
$\mathbb{T}_\Gamma = k[\![x^\gamma \mid \gamma \in \Gamma]\!]$:
\begin{equation}\label{eq:ks-symplecto}
 K_\gamma = \exp\Bigl(
 \sum_{n=1}^{\infty} \frac{(-1)^{n(|\gamma|+1)}}{n^2}
 \Omega(\gamma) \, x^{n\gamma}
 \Bigr) \in \Aut(\mathbb{T}_\Gamma).
\end{equation}
(Here the exponent acts via the Poisson bracket on $\mathbb{T}_\Gamma$
determined by the Euler form.)

For a half-plane $\ell \subset \mathbb{C}$ with $Z_t(\Gamma_+) \subset
\ell$, define the \emph{KS product}
\begin{equation}\label{eq:ks-product}
 \mathbf{A}_\ell(t)
 = \prod_{\gamma \in \Gamma_+,\, Z_t(\gamma) \in \ell}^{\curvearrowleft}
 K_\gamma^{\Omega(\gamma; t)}
\end{equation}
where $\curvearrowleft$ denotes the clockwise ordering by phase
$\Arg Z_t(\gamma)$.

\begin{theorem}[Kontsevich--Soibelman, 2008]
\label{thm:ks-wcf}
The ordered product $\mathbf{A}_\ell(t)$ is independent of $t$:
\begin{equation}\label{eq:ks-invariance}
 \mathbf{A}_\ell(t_+) = \mathbf{A}_\ell(t_-)
\end{equation}
for any two values $t_\pm$ on opposite sides of a wall of marginal
stability.
\end{theorem}

The content of the formula is that, although the individual factors
$K_\gamma^{\Omega(\gamma; t)}$ jump as $t$ crosses a wall, the product
(in the appropriate ordering) remains constant.

\subsection{Reinterpretation via MC gauge equivalence}
\label{ssec:ks-mc-gauge}

We now give the gauge-theoretic reinterpretation. The dictionary is:

\medskip
\begin{center}
\renewcommand{\arraystretch}{1.3}
\begin{tabular}{lll}
 \textbf{KS formalism} & & \textbf{MC/bar-cobar formalism} \\
 \hline
 Charge lattice $\Gamma$ & $\longleftrightarrow$ &
 Root lattice of $\frakg_X$ \\
 BPS index $\Omega(\gamma; t)$ & $\longleftrightarrow$ &
 Root multiplicity $\mult(\gamma)$ at moduli $t$ \\
 KS factor $K_\gamma$ & $\longleftrightarrow$ &
 Factor in bar-complex Euler product \\
 Ordered product $\mathbf{A}_\ell(t)$ & $\longleftrightarrow$ &
 Euler product $\Phi_X(x)$ \\
 Wall-crossing & $\longleftrightarrow$ &
 Gauge transformation of $\Theta_A(t)$ \\
 KS invariance & $\longleftrightarrow$ &
 Gauge invariance of $\Phi_X$ (Prop.~\ref{prop:euler-gauge-inv}) \\
\end{tabular}
\end{center}
\medskip

\begin{construction}[The wall-crossing gauge element]
\label{constr:wc-gauge}
Let $t_\pm \in \cM_K(X)$ be moduli on opposite sides of a wall
$\cW(\gamma_1, \gamma_2)$. The MC elements $\Theta_A(t_+)$ and
$\Theta_A(t_-)$ are related by a gauge transformation
\begin{equation}\label{eq:wc-gauge-element}
 \Theta_A(t_+) = e^{\alpha_{12}} \cdot \Theta_A(t_-)
\end{equation}
where the gauge element $\alpha_{12} \in (\frakg^{\mathrm{mod}}_A)^0$
is determined by the BPS states that bind or decay at the wall:
\begin{equation}\label{eq:alpha-formula}
 \alpha_{12}
 = \sum_{\substack{n_1, n_2 \geq 1 \\
 n_1 \gamma_1 + n_2 \gamma_2 \in \Gamma_+}}
 c(n_1, n_2; \Omega(\gamma_1), \Omega(\gamma_2))
 \, e_{n_1 \gamma_1 + n_2 \gamma_2}
\end{equation}
where $e_\gamma \in (\frakg^{\mathrm{mod}}_A)^0_\gamma$ are the
degree-zero Lie algebra generators of charge $\gamma$, and the
coefficients $c(n_1, n_2; \Omega_1, \Omega_2)$ are determined by
the BCH expansion of the KS factorization
$K_{\gamma_1}^{\Omega_1} K_{\gamma_2}^{\Omega_2}
= K_{\gamma_2}^{\Omega_2'} K_{\gamma_1+\gamma_2}^{\Omega_{12}'}
\cdots K_{\gamma_1}^{\Omega_1'}$.
\end{construction}

The key point: the KS wall-crossing formula \emph{is} the statement that
$\Theta_A(t_+)$ and $\Theta_A(t_-)$ are gauge equivalent. The ordered
product $\mathbf{A}_\ell$ is the gauge-invariant datum: it depends only
on the equivalence class $[\Theta_A] \in \overline{\MC}(\frakg^{\mathrm{mod}}_A)$.

\begin{conjecture}[Wall-crossing = MC gauge equivalence]
\label{conj:wc-mc-gauge}
Assume CY-A \textup{(}the CY-to-chiral functor at $d = 3$, proved for toric CY3, open in general\textup{)}.
Let $t_+, t_-$ be on opposite sides of a wall of marginal stability.
Then:
\begin{enumerate}[label=(\roman*)]
\item $\Theta_A(t_+)$ and $\Theta_A(t_-)$ lie in the same gauge orbit of
 $\MC(\frakg^{\mathrm{mod}}_A)$, i.e.,
 $[\Theta_A(t_+)] = [\Theta_A(t_-)]$ in
 $\overline{\MC}(\frakg^{\mathrm{mod}}_A)$.

\item The gauge transformation is $e^{\alpha_{12}}$ as in
 Construction~\ref{constr:wc-gauge}, and its BCH expansion reproduces
 the KS wall-crossing formula.

\item The bar-complex Euler product $\Phi_X$ is invariant:
 $\Phi_X(t_+) = \Phi_X(t_-)$. This is equivalent to the invariance of
 the denominator identity of $\frakg_X$.
\end{enumerate}
\end{conjecture}

\begin{remark}[Strategy \textup{(}conditional on CY-A\textup{)}]
Part (i) follows from the deformation-theoretic meaning of
$\overline{\MC}$: the chiral algebra $A_\cC$ is defined intrinsically
(it depends on $\cC$, not on a stability condition), so the MC moduli
class $[\Theta_A]$ is intrinsic. Different choices of stability condition
$\sigma$ give different \emph{representatives} $\Theta_A(\sigma)$ of
the same class.

Part (ii) is an explicit computation: the BCH formula for the gauge
action~\eqref{eq:gauge-action} applied to the $\Linf$-structure of
$\frakg^{\mathrm{mod}}_A$ reproduces the KS factorization algorithm.
The key input is that the $\Linf$-brackets $\ell_n$ encode the
boundary strata of $\overline{\cM}_{g,n}$, and the KS scattering
diagram is the tropical limit of this boundary stratification.

Part (iii) follows from Proposition~\ref{prop:euler-gauge-inv}: the
Euler product is a function on $\overline{\MC}$, not on $\MC$.
\end{remark}

\subsection{The scattering diagram as tropical MC data}
\label{ssec:scattering}

The Kontsevich--Soibelman \emph{scattering diagram} (also called the
``wall structure'' in the Gross--Siebert programme) admits a direct
interpretation in the $\Linf$ language. A scattering diagram
$\mathfrak{D}$ in $\Gamma_\mathbb{R} = \Gamma \otimes \mathbb{R}$
consists of walls (codimension-one cones) decorated by automorphisms of
$\mathbb{T}_\Gamma$. The consistency condition (that the
monodromy around any joint is trivial) is equivalent to the MC
equation for a \emph{tropical} $\Linf$-algebra obtained from
$\frakg^{\mathrm{mod}}_A$ by tropicalization (replacing moduli of curves
by their dual intersection complexes).

More precisely:

\begin{enumerate}[label=(\arabic*)]
\item Each wall in $\mathfrak{D}$ labelled by $\gamma$ with attached
 automorphism $K_\gamma^{\Omega(\gamma)}$ corresponds to a summand
 $\Omega(\gamma) \cdot \Theta_\gamma^{\mathrm{prim}} \, x^\gamma$ in
 the MC element $\Theta_A$.

\item The consistency condition (trivial monodromy around codimension-two
 joints) is the MC equation: $d\Theta + \frac{1}{2}[\Theta, \Theta]
 + \cdots = 0$. The higher $\Linf$-terms $\ell_n$ correspond to
 higher-codimension joints where $n$ walls meet.

\item Changing the half-plane $\ell$ (i.e., rotating the central charge
 ordering) corresponds to a gauge transformation of $\Theta_A$. The
 independence of $\mathbf{A}_\ell$ from $\ell$ is gauge invariance.
\end{enumerate}

% ======================================================================
\section{The attractor mechanism as canonical gauge}
\label{sec:attractor}
% ======================================================================

\subsection{Attractor flow and split BPS states}
\label{ssec:attractor-flow}

The \emph{attractor mechanism} (Ferrara--Kallosh--Strominger, 1995) states
that in $\cN = 2$ supergravity compactified on $X$, the moduli $t$ flow
along gradient lines towards \emph{attractor points} $t_*(\gamma)$
determined by the charge $\gamma$:
\begin{equation}\label{eq:attractor}
 Z_{t_*(\gamma)}(\gamma') \in \mathbb{R} \cdot Z_{t_*(\gamma)}(\gamma)
 \quad \text{for all } \gamma' \text{ with } \langle \gamma, \gamma'
 \rangle \neq 0.
\end{equation}
At the attractor point, the central charges of all constituents of a
potential bound state of charge $\gamma$ are aligned. The BPS index at
the attractor point is the \emph{single-centered} (or ``intrinsic'')
invariant $\Omega_*(\gamma) = \Omega(\gamma; t_*(\gamma))$.

\subsection{MC interpretation: canonical gauge}
\label{ssec:canonical-gauge}

In the MC language, the attractor point $t_*$ determines a distinguished
representative of the MC equivalence class:

\begin{definition}[Attractor MC element]
\label{def:attractor-mc}
The \emph{attractor MC element} is
\begin{equation}\label{eq:theta-attractor}
 \Theta_A^* = \sum_{\gamma \in \Gamma_+}
 \Omega_*(\gamma) \cdot \Theta_\gamma^{\mathrm{prim}} \, x^\gamma
 \in \MC(\frakg^{\mathrm{mod}}_A)
\end{equation}
where $\Omega_*(\gamma)$ is the single-centered (attractor) BPS index.
\end{definition}

This is the ``split'' gauge: the BKM root datum is diagonal in the sense
that each root space $(\frakg_X)_\gamma$ has dimension $|\Omega_*(\gamma)|$,
with no mixing between different charges.

At a generic point $t \neq t_*$, the MC element $\Theta_A(t)$ is obtained
from $\Theta_A^*$ by a gauge transformation that ``binds'' the
single-centered states into multi-centered bound states:
\begin{equation}\label{eq:binding-gauge}
 \Theta_A(t) = e^{\alpha(t, t_*)} \cdot \Theta_A^*
\end{equation}
where $\alpha(t, t_*) \in (\frakg^{\mathrm{mod}}_A)^0$ encodes the binding
data. The gauge element $\alpha(t, t_*)$ is the ``attractor flow tree''
datum: it is determined by the split attractor flow of Denef (2000).

\begin{remark}[The attractor as basepoint]
The attractor MC element $\Theta_A^*$ plays the role of a canonical
basepoint in $\MC(\frakg^{\mathrm{mod}}_A)$. The gauge orbit
$\exp((\frakg^{\mathrm{mod}}_A)^0) \cdot \Theta_A^*$ is the space of
all stability conditions compatible with the same intrinsic BPS spectrum.
Different attractor points (for different total charges $\gamma$) give
different canonical representatives, but they all lie in the same gauge
orbit.
\end{remark}

\subsection{Root multiplicities and gauge reshuffling}
\label{ssec:root-reshuffle}

At the attractor point, the root multiplicity function is
$\mult_*(\gamma) = |\Omega_*(\gamma)|$. Away from the attractor, the
gauge transformation $e^{\alpha(t, t_*)}$ reshuffles root multiplicities:
\begin{itemize}
\item Bound states appear: a composite charge $\gamma = \gamma_1 +
 \gamma_2$ acquires multiplicity from the binding of $\gamma_1$ and
 $\gamma_2$. In the $\Linf$-language, this is the bracket
 $\ell_2(\Theta_{\gamma_1}, \Theta_{\gamma_2})$ contributing to
 $\Theta_\gamma$.

\item Bound states decay: crossing a wall in the opposite direction
 reduces $\mult(\gamma)$ as the binding contribution is removed.
\end{itemize}

The total ``mass,'' measured by the bar-complex Euler product
$\Phi_X$, is invariant. What changes is the distribution of mass
among the root spaces. This is the gauge freedom in the MC
moduli space.

% ======================================================================
\section{The BKM denominator identity as gauge-invariant}
\label{sec:bkm-gauge-inv}
% ======================================================================

\subsection{The denominator identity for $K3 \times E$}
\label{ssec:denom-k3e}

For $X = K3 \times E$, the BKM superalgebra is $\frakg_{\Delta_5}$
and the denominator identity (product form) is
\begin{equation}\label{eq:denom-k3e}
 \frac{1}{64} \Delta_5
 = e^{\pi i(z_1 + z_2 + z_3)}
 \prod_{\substack{n,l,m \in \mathbb{Z} \\ (n,l,m) > 0}}
 (1 - e^{2\pi i(nz_1 + lz_2 + mz_3)})^{f(nm,l)}
\end{equation}
where $f(n,l)$ are the Fourier coefficients of $\phi_{0,1}$.

In the gauge-equivalence interpretation: the Igusa cusp form $\Delta_5$ is
the gauge-invariant datum. The root multiplicities $f(nm,l)$ appearing
in the product depend on the gauge (= choice of stability condition =
choice of fundamental domain in the K\"ahler moduli), but the product
itself is a Siegel modular form of weight 5 for $\mathrm{Sp}_4(\mathbb{Z})$
; it does not depend on the gauge.

\subsection{The denominator identity for toric CY3}
\label{ssec:denom-toric}

For a toric CY3 $X$, the denominator identity is the DT partition
function $Z^{\DT}(X)$. The DT/GW correspondence (MNOP conjecture, proved
in toric cases by Maulik--Nekrasov--Okounkov--Pandharipande) expresses
$Z^{\DT}$ in terms of topological string amplitudes.

The wall-crossing formula of KS is the statement that $Z^{\DT}$ is
independent of the stability condition. In the MC language: $Z^{\DT}$
depends only on the gauge equivalence class $[\Theta_A]$, not on the
representative.

\subsection{The general statement}
\label{ssec:general-euler-inv}

\begin{conjecture}[Denominator identity = gauge-invariant MC datum]
\label{conj:denom-gauge-inv}
For any CY3 threefold $X$ with quantum vertex chiral group $G(X)$,
the denominator identity of the BKM superalgebra $\frakg_X$ is the unique
(up to scalar) gauge-invariant formal power series on the MC moduli space
$\overline{\MC}(\frakg^{\mathrm{mod}}_A)$. Its automorphic properties
(modularity under $\Aut(\Gamma)$) descend from the $E_2$-structure of
$G(X)$.
\end{conjecture}

% ======================================================================
\section{The shadow obstruction tower and finite-order discontinuities}
\label{sec:shadow-tower-wc}
% ======================================================================

\subsection{Shadow obstruction tower truncations}
\label{ssec:shadow-trunc}

The shadow obstruction tower of Volume~I decomposes the MC element
$\Theta_A$ by degree:
\begin{equation}\label{eq:shadow-tower}
 \Theta_A = \Theta^{\leq 2}_A + \Theta^{(3)}_A + \Theta^{(4)}_A + \cdots
\end{equation}
where $\Theta^{\leq r}_A$ is the truncation to degrees $\leq r$. The
components are:
\begin{itemize}
\item $\Theta^{\leq 2}_A = \kappa_{\mathrm{ch}}$: the modular characteristic (= Weyl
 vector contribution in the BKM language). This is the ``tree-level''
 datum.
\item $\Theta^{(3)}_A = C$: the cubic shadow (first imaginary root
 corrections).
\item $\Theta^{(4)}_A = Q$: the quartic shadow (higher corrections).
\item The full tower $\Theta_A$: the complete automorphic correction.
\end{itemize}

\subsection{Truncations are not gauge-invariant}
\label{ssec:trunc-not-gauge-inv}

\begin{proposition}[Shadow obstruction tower discontinuity at walls]
\label{prop:shadow-jump}
For $r < \infty$, the truncation $\Theta^{\leq r}_A(t)$ may be
discontinuous at walls of marginal stability.
\end{proposition}

\begin{proof}[Argument]
The gauge element $\alpha_{12}$ of Construction~\ref{constr:wc-gauge}
has components at all degrees. The gauge action
$e^{\alpha_{12}} \cdot \Theta_A$ mixes degrees: the degree-$r$ component
of $e^{\alpha_{12}} \cdot \Theta_A$ depends on degrees $\leq r$ of
$\alpha_{12}$ and degrees $\leq r$ of $\Theta_A$, but also on cross-terms.
Truncating at degree $r$ destroys the gauge covariance.

Concretely: at a wall $\cW(\gamma_1, \gamma_2)$ with
$|\gamma_1| + |\gamma_2| = |\gamma|$, the gauge transformation shifts
$\Theta^{(|\gamma|)}_A$ by a term proportional to $[\alpha^{(|\gamma_1|)},
\Theta^{(|\gamma_2|)}_A]$. If this bracket is nonzero (which happens when
the BPS states genuinely bind or decay), then $\Theta^{\leq |\gamma|-1}_A$
jumps.
\end{proof}

\begin{example}[Conifold]
For the resolved conifold $X = \cO(-1) \oplus \cO(-1) \to \mathbb{P}^1$,
there is a single wall at $\Arg Z(\gamma_1) = \Arg Z(\gamma_2)$ where
$\gamma_1 = [\cO_{\mathbb{P}^1}]$ and $\gamma_2 = [\cO_p]$. On one
side of the wall, $\Omega(\gamma_1 + \gamma_2) = 1$ (the bound state
exists); on the other side, $\Omega(\gamma_1 + \gamma_2) = 0$. The
degree-2 truncation $\Theta^{\leq 2}_A$ jumps: on one side, it includes
the $\gamma_1 + \gamma_2$ contribution; on the other, it does not. The
full MC element $\Theta_A$ is continuous (gauge-equivalent on both sides).
\end{example}

\subsection{Physical interpretation: the renormalization group}
\label{ssec:rg-interpretation}

The shadow obstruction tower truncation has a physical interpretation as an
\emph{energy scale cutoff}. The degree-$r$ truncation $\Theta^{\leq r}_A$
captures BPS states with at most $r$ constituents. Wall-crossing at a
given wall can create or destroy states with a specific number of
constituents; the finite-order truncation sees this creation/destruction
as a discontinuity.

The full MC element $\Theta_A$ is the ``UV completion'': it includes all
degrees, and the gauge equivalence (= wall-crossing invariance) is restored.
This parallels the general principle that the full renormalization group
flow is scheme-independent, while finite-order truncations are
scheme-dependent.

% ======================================================================
\section{Relation to the BKM root datum}
\label{sec:bkm-root-wc}
% ======================================================================

\subsection{The root datum at different chambers}
\label{ssec:root-chambers}

The K\"ahler moduli space $\cM_K(X)$ is divided into chambers by walls
of marginal stability. In each chamber, the BPS spectrum
$\{\Omega(\gamma; t)\}_{\gamma \in \Gamma}$ is locally constant, and
determines a BKM root datum:
\[
 \cR_t(X) = (\Gamma, \Delta^{\mathrm{re}}_t, \Delta^{\mathrm{im}}_t,
 W_t, \rho_t, \mult_t).
\]
The root datum changes across walls: imaginary root multiplicities
$\mult_t(\gamma)$ jump. However:

\begin{proposition}
\label{prop:root-data-gauge-equiv}
The BKM root data $\cR_{t_+}(X)$ and $\cR_{t_-}(X)$ on opposite sides
of a wall are related by a ``gauge transformation of root data''
: a reshuffling of imaginary root multiplicities that preserves the
denominator identity $\Phi_X$.
\end{proposition}

This is the root-system manifestation of MC gauge equivalence: different
gauges give different root multiplicities, but the same denominator
identity (= the same automorphic form).

\subsection{The real roots are gauge-invariant}
\label{ssec:real-roots-inv}

The \emph{real} roots $\Delta^{\mathrm{re}}$ (those with
$(\alpha, \alpha) > 0$ and $\mult(\alpha) = 1$) do \emph{not} jump
across walls. In the MC language: the real roots correspond to the
degree-2 data $\kappa_{\mathrm{ch}} = \Theta^{\leq 2}_A$ at a ``sufficiently deep''
attractor point where no further binding is possible. Real roots are the
``hard'' generators; imaginary roots are the ``soft'' (gauge-dependent)
generators.

For $K3 \times E$: the three real simple roots
$\delta_1, \delta_2, \delta_3$ with Gram matrix
$\bigl(\begin{smallmatrix} 2 & -2 & -2 \\ -2 & 2 & -2 \\ -2 & -2 & 2
\end{smallmatrix}\bigr)$
are gauge-invariant. The imaginary root multiplicities $f(nm, l)$,
the Fourier coefficients of $\phi_{0,1}$, are gauge-dependent
(they depend on the choice of fundamental domain for the Weyl group
$W^{(2)}(\Lambda^{2,1}_{II})$), but the sum formula (denominator
identity $\Delta_5$) is gauge-invariant.

% ======================================================================
\section{Structural analogies and outlook}
\label{sec:wc-outlook}
% ======================================================================

\subsection{Analogy with Chern--Simons theory}
\label{ssec:cs-analogy}

The MC gauge equivalence interpretation of wall-crossing has a close
analogy with Chern--Simons theory on a 3-manifold $M$. There:
\begin{itemize}
\item The dgla is $\Omega^\bullet(M, \frakg)$, the dg Lie algebra of
 $\frakg$-valued differential forms.
\item MC elements are flat connections:
 $dA + \frac{1}{2}[A, A] = 0$.
\item Gauge equivalence is the usual gauge equivalence of flat connections.
\item The gauge-invariant datum is the Chern--Simons invariant
 $\mathrm{CS}(A)$, a function on the MC moduli space
 $\overline{\MC}(\Omega^\bullet(M, \frakg))$.
\end{itemize}

In our setting:
\begin{itemize}
\item The ``dgla'' is $\frakg^{\mathrm{mod}}_A$, the modular convolution
 algebra.
\item ``Flat connections'' are MC elements $\Theta_A$.
\item ``Gauge equivalence'' is wall-crossing.
\item The ``Chern--Simons invariant'' is the denominator identity
 $\Phi_X$ (or, in the K3$\times$E case, the Igusa cusp form $\Delta_5$).
\end{itemize}

This analogy can be made precise: the modular convolution algebra
$\frakg^{\mathrm{mod}}_A$ is a \emph{factorization} (= chiral) analogue
of $\Omega^\bullet(M, \frakg)$, where the 3-manifold $M$ is replaced by
the moduli of curves $\overline{\cM}_{g,n}$.

\subsection{Analogy with Stokes phenomena}
\label{ssec:stokes}

There is also an analogy with Stokes phenomena for irregular connections.
The ``Stokes data'' of an irregular connection on a punctured disk,
the Stokes matrices, which jump across Stokes rays, is gauge-equivalent
to the ``formal data'' (the formal monodromy and the irregular type).
The KS wall-crossing formula is the ``Stokes phenomenon'' for the MC
element $\Theta_A$: the BPS degeneracies $\Omega(\gamma; t)$ are the
Stokes multipliers, the walls of marginal stability are the Stokes rays,
and the gauge-invariant datum $\Phi_X$ is the ``connection class.''

This analogy was made precise by Bridgeland (2019) for certain classes
of CY3 categories, where the space of stability conditions
$\Stab(\cC)$ is identified with a space of irregular connections and
the wall-crossing formula becomes literally the isomonodromic
deformation equation.

\subsection{Implications for the programme}
\label{ssec:implications}

The MC gauge equivalence interpretation of wall-crossing has several
consequences for the quantum vertex chiral group programme:

\begin{enumerate}[label=(\arabic*)]
\item \textbf{Intrinsic definition of $G(X)$}: The quantum vertex chiral
 group $G(X)$ should be defined not by a specific BKM root datum
 (which is gauge-dependent), but by the MC moduli class
 $[\Theta_A] \in \overline{\MC}(\frakg^{\mathrm{mod}}_A)$. Different
 stability conditions give different presentations of the same
 algebraic object.

\item \textbf{The denominator identity as primary invariant}: The
 automorphic form (Igusa cusp form, DT partition function) is the
 primary invariant of $G(X)$: it is the unique gauge-invariant
 formal power series on the MC moduli space. The root multiplicities
 are gauge-dependent secondary data.

\item \textbf{Shadow obstruction tower and renormalization}: The shadow obstruction tower
 truncations $\Theta^{\leq r}_A$ provide ``effective'' descriptions of
 $G(X)$ at finite degree, analogous to effective field theories at a
 given energy scale. Wall-crossing discontinuities at finite order are
 the analogue of scheme dependence in perturbative QFT.

\item \textbf{The eight quantum vertex chiral groups}: The eight
 CY3 geometries $X_N = (S \times E)/(\mathbb{Z}/N\mathbb{Z})$ of
 Conjecture~\ref*{conj:eight-qvcg} should each give an MC moduli class,
 and the eight dd-modular forms are the gauge-invariant denominators.
 Wall-crossing \emph{between} different $X_N$ (i.e., deforming the
 K3$\times$E geometry) would correspond to a \emph{higher} gauge
 transformation connecting different MC moduli spaces.
\end{enumerate}

\subsection{Open questions}
\label{ssec:wc-open}

\begin{enumerate}[label=(Q\arabic*)]
\item Is the gauge element $\alpha_{12}$ of
 Construction~\ref{constr:wc-gauge} determined by the split attractor
 flow tree data of Denef? Can the flow tree formula be derived from
 the $\Linf$-structure of $\frakg^{\mathrm{mod}}_A$?

\item The KS wall-crossing formula applies to ``motivic''
 (= categorified) DT invariants. Is there a \emph{categorified} MC
 gauge equivalence, where $\frakg^{\mathrm{mod}}_A$ is replaced by an
 $\Linf$-\emph{category} and gauge equivalence is a natural
 transformation?

\item For compact CY3s (e.g., the quintic), the BPS spectrum is
 largely unknown. Can the MC gauge equivalence framework constrain
 the spectrum via the requirement that $[\Theta_A]$ be well-defined?

\item What is the relationship between the MC gauge group
 $\exp((\frakg^{\mathrm{mod}}_A)^0)$ and the autoequivalence group
 $\Aut(\cC)$ of the CY category? Is every wall-crossing gauge
 transformation induced by an autoequivalence (e.g., a
 Fourier--Mukai transform)?

\item The shadow obstruction tower truncation $\Theta^{\leq r}_A$ jumps at walls.
 Is there a ``regularized'' truncation (analogous to a renormalization
 scheme) that is continuous across walls, at the cost of introducing
 scheme-dependent counterterms?
\end{enumerate}

% ======================================================================
\section{Wall-crossing as MC gauge: swarm verification}
\label{sec:wc-swarm}
% ======================================================================

The identification of wall-crossing with MC gauge equivalence
has been verified computationally (73~tests in
\texttt{compute/lib/wallcrossing\_gauge\_engine.py}, 124~tests
in \texttt{compute/lib/conifold\_bar\_complex.py}). We record
the key findings.

\subsection{Pentagon identity from the Jacobi identity}

The KS pentagon identity (consistency of the scattering diagram
around a codimension-2 joint where five walls meet) is a
consequence of $D^2 = 0$ in the convolution algebra. More
precisely:

\begin{enumerate}[label=(\roman*)]
 \item The $\Linf$ Jacobi identity $\sum_{i+j = n+1}
 \ell_i \circ \ell_j = 0$ (which is $D^2 = 0$ unpacked) implies
 the consistency of scattering at every joint.

 \item The pentagon identity for five BPS rays in the upper
 half-plane is the Jacobi identity at degree~3, projected onto
 the charge lattice. This is verified explicitly for the conifold.

 \item Higher consistency conditions (hexagons, etc.) arise from
 higher $\Linf$-brackets at degrees $\geq 4$.
\end{enumerate}

\subsection{Reineke vs fermionic conventions}

Two conventions for the KS automorphism $K_\gamma$ appear in the
literature:

\begin{itemize}
 \item \textbf{Reineke (rational DT):}
 $K_\gamma = \exp(-\sum_{k \geq 1} e_{k\gamma}/k^2)$.
 This is the motivic Hall algebra convention, natural for
 rational DT invariants $\overline{\Omega}(\gamma)$.

 \item \textbf{Fermionic (integer DT):}
 $K_\gamma = \prod_{k \geq 1} (1 - (-1)^k e_{k\gamma})^{k\, \Omega(k\gamma)}$.
 This is the physical convention, natural for integer BPS
 degeneracies $\Omega(\gamma)$.
\end{itemize}

\noindent
The two conventions produce different gauge elements $\alpha_{12}$
but the same gauge-equivalence class $[\Theta_A]$. The compute
layer verifies this explicitly for the conifold: both conventions
yield the same $\overline{\MC}$ class.

\subsection{Conifold: gauge terminates at order 1}

For the resolved conifold with two BPS hypermultiplets of charges
$\gamma_1, \gamma_2$ (both with $\Omega = 1$ and
$\langle \gamma_1, \gamma_2 \rangle = 1$), the gauge
transformation terminates:

\[
 \alpha_{12} = e_{\gamma_1 + \gamma_2}, \qquad
 [\alpha_{12}, \alpha_{12}] = 0.
\]

\noindent
There is a single bound state of charge $\gamma_1 + \gamma_2$
with $\Omega(\gamma_1 + \gamma_2) = 1$, and no higher-order
BCH corrections. On the bar-complex side, $B(A_{\mathrm{con}})$
has exactly three nonzero root-graded pieces (charges $\gamma_1$,
$\gamma_2$, $\gamma_1 + \gamma_2$), each one-dimensional.
The conifold bar complex is a three-term chain complex whose
Euler characteristic reproduces the DT partition function
$M(q)^2 (1 - Q) = M(q)^2 \sum_n (-Q)^n$.

\begin{remark}[Gauge termination and shadow depth]
The conifold has $r_{\max} = 3$ (class~$L$): the cubic shadow
$C$ is the last nonzero obstruction. Gauge termination at BCH
order~1 is consistent with the tower truncating at degree~3.
For CY3s with richer BPS spectra (class~$M$, infinite
tower), the gauge transformation involves arbitrarily high BCH
orders, matching the infinite degree tower.
\end{remark}

% ======================================================================
\section*{References}
\label{sec:wc-refs}
% ======================================================================

\begin{enumerate}[label={[\arabic*]}, leftmargin=*, itemsep=2pt]
\item \label{ref:ks08} M.~Kontsevich and Y.~Soibelman,
 \textit{Stability structures, motivic Donaldson--Thomas invariants
 and cluster transformations}, arXiv:0811.2435 (2008).

\item \label{ref:fks95} S.~Ferrara, R.~Kallosh, and A.~Strominger,
 \textit{$\cN = 2$ extremal black holes},
 Phys.\ Rev.\ D~\textbf{52} (1995) R5412.

\item \label{ref:denef00} F.~Denef,
 \textit{Supergravity flows and D-brane stability},
 JHEP \textbf{08} (2000) 050.

\item \label{ref:bridgeland19} T.~Bridgeland,
 \textit{Riemann--Hilbert problems from Donaldson--Thomas theory},
 Invent.\ Math.\ \textbf{216} (2019) 69--124.

\item \label{ref:gps10} M.~Gross, R.~Pandharipande, and B.~Siebert,
 \textit{The tropical vertex},
 Duke Math.\ J.\ \textbf{153} (2010) 297--362.

\item \label{ref:mnop06} D.~Maulik, N.~Nekrasov, A.~Okounkov, and
 R.~Pandharipande,
 \textit{Gromov--Witten theory and Donaldson--Thomas theory, I},
 Compos.\ Math.\ \textbf{142} (2006) 1263--1285.

\item \label{ref:sv13} O.~Schiffmann and E.~Vasserot,
 \textit{Cherednik algebras, W-algebras and the equivariant
 cohomology of the moduli space of instantons on $\mathbb{A}^2$},
 Publ.\ Math.\ IH\'ES \textbf{118} (2013) 213--342.

\item \label{ref:rsyz} M.~Rapcak, Y.~Soibelman, Y.~Yang, and G.~Zhao,
 \textit{Cohomological Hall algebras, vertex algebras, and
 instantons},
 Comm.\ Math.\ Phys.\ \textbf{376} (2020) 1803--1873.

\item \label{ref:joyce21} D.~Joyce,
 \textit{Enumerative invariants and wall-crossing formulae in
 abelian categories},
 arXiv:\allowbreak 2111.04694 (2021).

\item \label{ref:gmn10} D.~Gaiotto, G.~W.~Moore, and A.~Neitzke,
 \textit{Wall-crossing, Hitchin systems, and the WKB approximation},
 Adv.\ Math.\ \textbf{234} (2013) 239--403.
\end{enumerate}

\end{document}
